%
%  Vorlage/Template fuer #Proseminar
%
%  Created by Claudia Pittel und Yunus Erdemir in 2018, updated in 2021.
%  Copyright (c) 2021 . All rights reserved.
%
\documentclass[12pt,toc=bib,toc=listof]{scrreprt}
\usepackage[ngerman]{babel} 
\usepackage[utf8]{inputenc}
\usepackage[T1]{fontenc}
% Standardschriftart Latin Modern
\usepackage[scaled]{helvet}
\renewcommand\familydefault{\sfdefault}
% Arial, so in etwa
%\usepackage{newtxmath,newtxtext}
%Times New Roman, so in etwa
% Zitate durchgehend nummerieren
\usepackage{chngcntr}
\counterwithout{footnote}{chapter}
%Zeilenabstand
\usepackage{setspace}
\linespread{1.3}

% DIN-A-4, Ränder (oben, links, rechts, unten), Abstand Kopfzeile und Text
\usepackage{geometry}
\geometry{a4paper, top=25mm, left=30mm, right=20mm, bottom=25mm, headsep=12mm}

\usepackage[backend=bibtex, style=authortitle-ibid, url=false, isbn=false, doi=false]{biblatex}
\addbibresource{literature.bib}

\usepackage{hyperref}
\hypersetup{
  ,colorlinks=true
  ,linkcolor=blue
  ,citecolor=blue
  ,filecolor=blue
  ,urlcolor=blue
  }

%%%%%%%%%%%%%%%%%%%%%%%%%%%%%%%%%%%%% % (fold)
% Vom Studierenden zu aendernde Werte
\newcommand{\topic}{XaaS - Anything as a Service}

\newcommand{\studentnameA}{Suphi Pembe}
\newcommand{\studentidA}{207617}
\newcommand{\studentpartA}{SEITEN oder KAPITEL VON BIS 1}

\newcommand{\studentnameB}{Andreas Würzer}
\newcommand{\studentidB}{207258}
\newcommand{\studentpartB}{SEITEN oder KAPITEL VON BIS 2}

\newcommand{\studentnameC}{VORNAME NACHNAME 3}
\newcommand{\studentidC}{MATRIKELNUMMER 3}
\newcommand{\studentpartC}{SEITEN oder KAPITEL VON BIS 3}

\newcommand{\semester}{Sommersemester 2022}
%
%%%%%%%%%%%%%%%%%%%%%%%%%%%%%%%%%%%%% % (end)

\usepackage{ifpdf}
\ifpdf
\usepackage[pdftex]{graphicx}
\else
\usepackage{graphicx}
\fi

\usepackage[headsepline,footsepline]{scrlayer-scrpage}
\pagestyle{scrheadings}
\clearpairofpagestyles
\ihead{Proseminar: \topic}
\ofoot{Seite \pagemark}
\renewcommand*{\chapterpagestyle}{scrheadings}
\renewcommand*{\chapterheadstartvskip}{}

% Logo
\titlehead{\flushright\includegraphics[scale=0.5]{HHN_Logo_D_oS_RGB_300.png}}
\subject{Proseminar (282136)}
\title{\topic}
\author{\studentnameA { ({\studentidA)}}, \\ \studentnameB { ({\studentidB)}}, \\ \studentnameC { ({\studentidC)}} }
\date {\semester}
%% Datum nie auf einen festen Wert setzen
\publishers{Vorgelegt bei Claudia Pittel}

%\pagestyle{headings}

% Zähler für Römische Nummerierung
\newcounter{savepage}

\begin{document}
\pagenumbering{roman} 
\selectlanguage{ngerman}
\sffamily
% serifenlose Schriftfamilie
\maketitle

\addchap{Management Summary} % (fold)
\label{sec:management_summary}

Hier sollte ziemlich genau bzw. maximal 1 Seite Text stehen (ziemlich genau  bedeutet, man sollte so nah wie möglich an 1 Seite herankommen).
Text für Test commit haha

% chapter management_summary (end)

\tableofcontents

\addchap{Abkürzungsverzeichnis} % (fold)
\label{sec:abkuerzungsverzeichnis}

\begin{description}
\item[GPU] Graphics-Processing-Unit oder Grafikkarte
\item[HPC] High-Performance-Computing
\end{description}

% chapter abkuerzungsverzeichnis (end)

\listoffigures
\listoftables

% \onehalfspacing

\newpage
% Zähler speichern
\setcounter{savepage}{\arabic{page}}
\pagenumbering{arabic}

\chapter{Einleitung} % (fold)
\label{sec:einleitung}

Einleitungstext mit Motivation, Ziel der Arbeit (d.h. Erläuterung der Forschungsfrage) und Beschreibung der Vorgehensweise bzw. Aufbau der Arbeit\footcite [vgl.] [S. 38]{Th17}

\section{Motivation} % (fold)
\label{sec:motivation}
%%%% Version 1 %%%%%%%%
Durch den aktuell anhaltende Halbleitermangel besteht ein Engpass an Ressourcen von die meisten Wirtschaftszweige betroffen sind.
Einer dieser Wirtschaftszweige ist die Produktion von GPUs (graphics processing unit. Diese werden für diverse Anwendung von Computern verwendet, im betrieblichen wie auch im privaten Bereich.
Primär in dieser Arbeit werden die Bereiche High-Performance-Computing (HPC) und Gaming haben.
Beide diese Bereiche benötigen GPU-Rechenleistung, welche im konventionell von einer lokal verbauten GPU zur Verfügung gestellt wird. 
Als Langfristige Lösung soll analysiert werden ob es möglich ist durch zentrale Services, welche GPU as a Service anbieten. 
Durch die zentralen Ressourcenteilung dem Mangel entgegenzuwirken mit einer alternative für den Bedarf zu schaffen.

% section motivation (end)

\section{Ziel der Arbeit} % (fold)
\label{sec:ziel_der_arbeit}

\ldots

% section ziel_der_arbeit (end)

\section{Vorgehensweise} % (fold)
\label{sec:vorgehensweise}

\ldots

% section vorgehensweise (end)
% chapter einleitung (end)

\chapter{Anything as a Service - Cloud Computing} % (fold)
\label{sec:Anything as a Service - Cloud Computing}

Zwischen den Gliederungspunkten sollten jeweils kurze Überleitungssätze stehen, damit man weiß, um was es inhaltlich in den folgenden Unterkapiteln geht.\footcite [Vgl.] []{GPL}

% chapter ersteskapitel (end)

\section{Definition} % (fold)
\label{sec:Definition}

Bei den Gliederungspunkten immer auf eine Ausgewogenheit achten, damit eine gleichmäßige Gliederung gefördert werden kann. Sofern Abbildungen (wie Abbildung 1: Beispielbild) verwendet werden, müssen diese auch inhaltlich im Text erwähnt und erläutert werden, sowie ein Abbildungsverzeichnis erstellt werden.\footcite [Vgl.] [] {hhnwin}

% section unterkapitel1 (end)

\section{Typische Servicemodelle} % (fold)
\label{sec:Typische Servicemodelle}

Untergliederungen nur in der Mehrzahl erstellen, d.h. nie 1 Unterkapitel alleine stehen lassen.
In gleicher Art und Weise wie Abbildungen dargestellt und beschriftet werden, verhält es sich mit Tabellen.

% section unterkapitel1 (end)

\subsection{IaaS: Infrastructure as a Service}
\label{sec:IaaS: Infrastructure as a Service}

Inhalt

\subsection{SaaS: Software as a Service}
\label{sec:SaaS: Software as a Service}

Inhalt

\subsection{PaaS: Plattform as a Service}
\label{sec:PaaS: Plattform as a Service}


\section{Vor- und Nachteile}
\label{Vor- und Nachteile}


\chapter{Knappheit von Grafikkarten} % (fold)
\label{sec:Knappheit von Grafikkarten}

%Sample Text%Der Hauptteil sollte nicht nur aus einem Oberkapitel bestehen, sondern aus mehreren, so dass die Hauptarbeit auch vom Umfang her im Hauptteil liegt. Dementsprechend sollten die Oberkapitel (erste Gliederungsebene) auch eigenständige Kapitelüberschriften besitzen.
%%%%%%%%%%%%%%% ROHSCHRIFT IDEE%%%%%%%%%%%%%%%%%%%%%%%%%%%%%%%%%%%
--------  ROHSCHRIFT UND IDEEN  --------\\
Die Knappheit von Grafikkarten, in Zukunft als GPU bezeichnet, hat den
aktuellen Markt durch neue Branchen die GPU-Leistung nutzen nachhaltig verändert.
In diesem Kapitel soll die Preisentwicklung der der GPUs betrachtet und im Zusammenhang 
gebracht werden mit dem Bedarf der Branchen. Dies wird aufzeigen, dass die Preisentwicklung durch
eine gleichzetige erhöhung des Bedarfs und fehlenden Verfügbarkeit geschehen ist.\\
%Paper finden das so eine Datenanalyse aufweißt
Ebenfalls wird erläutert wie diese Probleme konkret diese Situation beeinflussen und wie diese 
Langfristig zu bewerten sind.
%%%%%%%% Quellenanalyse nochmal checken!
%Paper Auswirkung Halbleitermangel
%Paper Wie ist Halbleitermangel
% https://ieeexplore.ieee.org/document/9568259
%BUCH Was ist Kryptomining
% https://link.springer.com/book/10.1007/978-3-030-97951-5
%Paper Problem an Kryptomining
% https://ieeexplore.ieee.org/document/9027372
%%%%%%%%%%%%%%%% ROHSCHRIFT IDEE%%%%%%%%%%%%%%%%%%%%%%%%%%%%%%%%%%%%%%%%%%%%


% chapter ersteskapitel (end)
%%% ggf. weitere Abschnitte

\section{Preisentwicklung}
\label{sec:Preisentwicklung}

Inhalt

\section{Ursache Halbleitermangel und KryptoMining}
\label{sec:Ursache Halbleitermangel und KryptoMining}

Inhalt

\chapter{Gaming as a Service}
\label{sec: Gaming as a Service}

Inhalt

\section{Funktionsweise}
\label{sec:Funktionsweise}

Inhalt

\section{Anbietervergleich}
\label{sec:Anbietervergleich}

Inhalt

\subsection{Voraussetzung}
\label{sec:Vorraussetzung}

Inhalt

\subsection{Angebot}
\label{sec:Angebot}

Inhalt

\subsection{Preis}
\label{sec:Preis}

Inhalt

\section{Hardwarevorraussetzung um Usability zu gewährleisten}
\label{sec:Hardwarevorraussetzung um Usability zu gewährleisten}

Inhalt

\chapter{GPU as a Service}
\label{sec:GPU as a Service}

Inhalt

\section{Funktionsweise}
\label{sec: Funktionsweise}

Inhalt

\section{Einsatzgebiete}
\label{sec: Einsatzgebiete}

Inhalte

\section{Vergleich eigene GPU und GPU in der Cloud}
\label{sec:Vergleich eigene GPU und GPU in der Cloud}

Inhalt

\chapter{Marktvorhersage}
\label{sec:Marktvorhersage}

Inhalt

\chapter{Fazit und Ausblick} % (fold)
\label{sec:fazit}

Kritische Begutachtung inklusive Zusammenfassung der Arbeit sowie eventuelle Zukunftsperspektiven zum Thema können hier im Fazit und im Ausblick eingebracht werden.\footcite [Vgl.] [] {HAN2016S30}

% chapter fazit (end)

\appendix
\newpage

\pagenumbering{roman}
% Zähler laden
\setcounter{page}{\thesavepage}

\addchap{Anhang} % (fold)
\label{sec:anhang}

Der Anhang soll den eigentlichen Hauptteil nicht ergänzen, sondern darüber hinaus weitere möglicherweise interessante Informationen liefern, die aber nicht zwangsläufig notwendig sind, um den Hauptinhalt zu verstehen

\newpage

% Quellen wird erst dargestellt sobald ein Zitat verwendet wird
\defbibheading{head}{\addchap{Quellenverzeichnis}}
\printbibliography[heading=head]

\vspace{2cm}

Hier müssen alle Quellenverweise zu finden sein – inklusive aller erforderlichen Angaben, alphabetisch sortiert. Eine Unterteilung in verschiedene Quellenarten ist grundsätzlich nicht notwendig, da die unterschiedlichen Quellenarten anhand der Angabe der bibliographischen Angaben zu erkennen ist. (d.h. beispielsweise keine Unterteilung zwischen „Printquellen“ und „Internetquellen“!)

\newpage

\addchap{Ehrenwörtliche Erklärung} % (fold)
\label{sec:erklaerung}

„Wir versichern, dass die vorliegende Arbeit von uns selbständig und ausschließlich unter Verwendung der angegebenen Quellen und Hilfsmittel angefertigt wurde. Alle Stellen, die wörtlich oder annähernd aus Veröffentlichungen entnommen sind, haben wir als solche kenntlich gemacht. Die Arbeit wurde bisher in gleicher oder ähnlicher Form, auch nicht in Teilen, keiner anderen Prüfungsbehörde vorgelegt und auch nicht veröffentlicht.“

\vspace{1cm}
\noindent
{\studentpartA} wurden von {\studentnameA} verfasst.
\\
{\studentpartB} wurden von {\studentnameB} verfasst.
\\
{\studentpartC} wurden von {\studentnameC} verfasst.

\vspace{3cm}
Ort, Datum \hfill Unterschrift

\vspace{2cm}
Ort, Datum \hfill Unterschrift

\vspace{2cm}
Ort, Datum \hfill Unterschrift


\end{document}
